% v2_ch8_delivery.tex — Chapter 8: Delivery: Browser, Engine, and Archive
% Video-to-Gaussian: An Agentic Pipeline for Volumetric Cultural Heritage
% Preservation (v2.0)
% Dr John O'Hare, DreamLab AI Consulting Ltd

\chapter{Delivery: Browser, Engine, and Archive}
\label{ch:delivery}

The reconstruction pipeline described in Chapters~\ref{ch:pipeline}--\ref{ch:scenegraph}
produces a single hierarchical \acrshort{usd} scene graph that encodes the full scene:
the trained Gaussian splat, the environment mesh, thirty-one per-object generative
meshes emplaced at surveyed positions, and associated catalogue metadata. From that
single artefact the \texttt{gaussian-toolkit} supports three distinct delivery paths
simultaneously---a compressed splat for real-time browser viewing, a mesh hierarchy
for game-engine and immersive workflows, and a preservation-grade \acrshort{usd}
archive for long-term custody. This chapter documents each path and justifies the
design decision to maintain all three from one master representation.

% ─────────────────────────────────────────────────────────────────────────────
\section{Compressed Splat Delivery (\texttt{.ksplat})}
\label{sec:delivery-ksplat}
% ─────────────────────────────────────────────────────────────────────────────

Raw \acrshort{3dgs} training produces a \texttt{.ply} point cloud in which each
Gaussian is stored as a 59-float record covering position, covariance, colour
harmonics, and opacity. For the gallery test scene this amounted to 1.18 million
Gaussians occupying approximately 267\,MB on disk---too large for low-latency
progressive streaming over a typical institution's public-facing web connection.

The \texttt{splat\_optimizer.py} module addresses this by wrapping the PlayCanvas
\texttt{splat-transform} tool \parencite{SplatTransform2025} in a Python driver
that applies four sequential operations before producing the final \acrshort{ksplat}
container.

\paragraph{Region-of-interest crop.}
Gaussians outside an axis-aligned bounding box defined either interactively through
the web \acrshort{usd} interface or by the scene-graph bounding volumes are
discarded. For the gallery scene this removed approximately 7\,\% of Gaussians
corresponding to floater artefacts above the ceiling plane and behind the capture
boundary.

\paragraph{Opacity and scale filtering.}
Gaussians with opacity below a configurable threshold (default 0.05) contribute
negligible visual energy and are pruned. Gaussians with scale in any principal axis
exceeding a configurable maximum (default 2\,m) are artefacts of under-constrained
geometry and are likewise removed. Together these two filters account for an
additional 6\,\% reduction.

\paragraph{Morton-order spatial sort.}
The surviving Gaussians are reordered according to their Morton (Z-order) curve
index computed from their quantised 3-D position. Morton ordering interleaves the
bits of the three spatial coordinates to produce a one-dimensional index that
preserves spatial locality: Gaussians that are near one another in 3-D space appear
consecutively in the sorted sequence. When the browser loads the file progressively,
each kilobyte of data decoded corresponds to a spatially coherent region of the
scene rather than a globally scattered sample. The practical effect is that the
viewer renders a coarse-but-complete version of the whole scene in the first few
hundred milliseconds of loading, then refines it as subsequent data arrives,
rather than revealing the scene piecemeal from one side to the other.

\begin{technote}
Morton ordering does not reduce file size; it redistributes data for streaming
locality. Compression of the quantised attributes is handled by the
\texttt{splat-transform} codec, which applies per-attribute range quantisation
followed by entropy coding. The two steps are independent and complementary.
\end{technote}

\paragraph{Compression.}
After sorting, \texttt{splat-transform} quantises positions to 16-bit integers
within the scene bounding box, reduces spherical-harmonic degree from three to
one (retaining the dominant view-dependent colour component), and encodes the
resulting attributes using a compact binary format. The resulting \acrshort{ksplat}
file for the 1.18-million-Gaussian gallery scene is 18\,MB---a 14.8$\times$ reduction
relative to the raw \texttt{.ply}.

The bandwidth-quality trade is straightforward: harmonic truncation introduces
minor view-dependent artefacts on highly specular surfaces (notably polished display
case glass) that are imperceptible on matte painted walls and organic sculptural
forms. For cultural heritage preview delivery this trade is almost always
favourable; institutions that require full spectral fidelity for technical
examination can serve the uncompressed \texttt{.ply} to specialist users on request.

% ─────────────────────────────────────────────────────────────────────────────
\section{Browser Viewing}
\label{sec:delivery-browser}
% ─────────────────────────────────────────────────────────────────────────────

The \acrshort{ksplat} file is served by the toolkit's built-in static file handler
and rendered in the browser via a \acrshort{webgl}/\acrshort{webgpu} Gaussian
splat renderer \parencite{Antimatter15Splat,Khronos2024WebGPU}. No plugin, extension,
or application installation is required; a visitor with a contemporary browser on
any operating system can view the full photorealistic scene interactively on
consumer hardware. The gallery reconstruction ran at 160\,FPS at 1080p on a
mid-range desktop \acrshort{gpu} and at a comfortable 45--55\,FPS on an integrated-graphics
laptop.

\begin{figure}[htbp]
  \centering
  \includegraphics[width=0.85\textwidth]{v2/webui.png}
  \caption{The \texttt{gaussian-toolkit} delivery and control web interface. The left
    panel provides scene controls and export triggers; the main viewport renders the
    \acrshort{ksplat} scene in real time. The info overlay at lower right displays
    the catalogue metadata associated with the currently selected object prim.}
  \label{fig:webui}
\end{figure}

\subsection{Object Picking and Metadata Surfacing}

The viewer is integrated with the \acrshort{usd} scene graph described in
Chapter~\ref{ch:scenegraph}. When a visitor clicks on the rendered scene the viewer
resolves the screen-space pick through a secondary depth-sorted instancing pass and
maps the hit back to the originating scene-graph prim. The prim carries the object's
catalogue metadata---accession number, artist, title, medium, date, condition
notes---as \acrshort{usd} custom attributes. On a successful pick the viewer's
overlay panel populates with this information without a round-trip to any external
database, since the \acrshort{usd} scene graph is self-contained and fully loaded
in the browser session.

This integration means that a visitor can click on an artwork in the
reconstruction and immediately receive its curatorial record, replicating the
experience of a digital label without requiring separate data services. The
picking resolution operates at prim granularity; future work could extend it to
Gaussian cluster granularity for sub-object selection.

\subsection{Viewer Option Comparison}

Two deployment configurations are supported: self-hosted \texttt{gsplat.js} and
the toolkit's integrated viewer. Table~\ref{tab:viewer-comparison} summarises
the relevant characteristics for an institution choosing between them.

\begin{table}[htbp]
  \centering
  \caption{Browser splat viewer options for institutional deployment.}
  \label{tab:viewer-comparison}
  \begin{tabular}{@{}L{3.2cm}L{2.8cm}L{3.8cm}@{}}
    \toprule
    \textbf{Feature}      & \textbf{Self-hosted gsplat.js} & \textbf{Toolkit integrated viewer} \\
    \midrule
    Max Gaussians tested  & $\sim$5\,M                     & $\sim$2\,M (\acrshort{ksplat} compressed) \\
    \acrshort{webgpu} support & Yes (fallback to \acrshort{webgl}) & Yes (\acrshort{webgl} primary) \\
    Mobile support        & Limited (iOS Safari constrained) & Reduced quality, functional \\
    USD prim picking      & Not included                   & Integrated (Chapter~\ref{ch:scenegraph}) \\
    Catalogue metadata overlay & Not included              & Integrated \\
    Licence               & MIT                            & MIT (DreamLab) \\
    Deployment complexity & Low (single JS bundle)         & Moderate (Docker service) \\
    \bottomrule
  \end{tabular}
\end{table}

Self-hosted \texttt{gsplat.js} is appropriate when the institution already manages
its own web infrastructure and wishes to embed viewing within an existing collection
management system. The toolkit viewer is appropriate when the institution requires
integrated metadata surfacing and object picking without additional bespoke
development. Both paths serve the same \acrshort{ksplat} file.

% ─────────────────────────────────────────────────────────────────────────────
\section{Engine and Immersive Delivery}
\label{sec:delivery-engine}
% ─────────────────────────────────────────────────────────────────────────────

The \acrshort{usd} scene graph exported by the pipeline is directly importable into
Unreal Engine 5.4 and later via Epic's native \acrshort{usd} importer and a
Gaussian splat rendering plugin \parencite{Pixar2019}. This import path preserves
full reconstruction fidelity: the 1.18-million-Gaussian splat is rendered natively
by the \acrshort{gpu} without conversion to a polygonal proxy, retaining all
view-dependent colour variation and sub-pixel surface detail.

\subsection{Variant-Set Architecture}

Within the exported \acrshort{usd} file each object prim carries a variant set with
two named variants: \texttt{Gaussian} and \texttt{Mesh}. Switching the variant set
at the \acrshort{usd} level substitutes the Gaussian representation with the
corresponding generative polygonal mesh for every object in the scene simultaneously.

This architecture has a practical consequence that simplifies immersive production
workflows considerably: a single asset pipeline serves both a photoreal splat viewer
for public browser access and a fully interactive \acrshort{pbr}-lit polygon scene
for \acrshort{vr} experiences and large-format display, without maintaining separate
files or repeating the reconstruction. A \acrshort{vr} developer activates the
\texttt{Mesh} variant to obtain collision geometry, physically-based material
assignment, and rigid-body physics; a web publisher activates the
\texttt{Gaussian} variant for photoreal browser delivery. The surveyed object
positions, catalogue metadata, and scene hierarchy are identical in both variants
because they are inherited from the common prim structure rather than stored per
variant.

\begin{keyfinding}
The \acrshort{usd} variant-set pattern resolves the central tension between splat
fidelity and mesh utility by delivering both from a single capture and a single
scene-graph artefact. Downstream consumers select the representation appropriate to
their application without any reprocessing.
\end{keyfinding}

\subsection{Mesh Variant for Collision and Large-Format Use}

The mesh variant is particularly important for three use cases that the Gaussian
variant cannot satisfy: collision detection in \acrshort{vr} locomotion systems,
where ray casting against Gaussian primitives is not yet standardised in any engine;
large-format physical installation, where tiled LED walls or projection mapping
require explicit polygon geometry for edge blending; and 3D printing of individual
objects, where watertight mesh topology is a hard requirement. In all three cases the
generative per-object meshes produced by the pipeline (Chapter~\ref{ch:objects})
provide an immediately usable foundation.

% ─────────────────────────────────────────────────────────────────────────────
\section{Archival Preservation}
\label{sec:delivery-archive}
% ─────────────────────────────────────────────────────────────────────────────

The canonical preservation artefact produced by the \texttt{gaussian-toolkit} is
the \acrshort{usd} scene graph itself. This single file encodes:

\begin{itemize}
  \item the full-resolution Gaussian splat (\texttt{.ply} referenced as a sublayer);
  \item the environment mesh produced by the selected extraction backend
        (\acrshort{tsdf} or \acrshort{milo});
  \item thirty-one per-object generative meshes at surveyed positions;
  \item camera poses from the \acrshort{sfm} reconstruction (for future
        re-training or quality audit);
  \item per-object catalogue metadata as custom prim attributes; and
  \item a \texttt{variant\_info} metadata block recording the software versions,
        training parameters, and extraction backend used.
\end{itemize}

The format is open and vendor-neutral, standardised by Pixar and adopted across the
visual effects, games, and automotive industries \parencite{Pixar2019}. Files written
today will remain readable by any compliant \acrshort{usd} implementation regardless
of whether the originating pipeline software continues to exist.

\subsection{Splat-as-Artefact Argument}

A recurring debate in digital preservation practice concerns whether a derived
representation---a mesh, a render, a downsampled image---is an adequate archival
substitute for the source data. For Gaussian splatting the answer is unambiguous:
the splat is the most information-rich representation produced by the pipeline.
Every mesh extracted from it is a lossy derivative that discards view-dependent
colour, opacity, and the fine-scale structure of surfaces that resisted meshing.
Once that information is absent from the archive it cannot be recovered.

Preserving the splat as the primary artefact therefore follows the same logic as
preserving camera RAW files alongside processed JPEGs: future operators will benefit
from improved meshing algorithms, improved compression schemes, and improved
rendering hardware that did not exist at capture time, but only if the source data
from which those improvements can be applied has been retained. A \acrshort{usd}
archive containing the full-resolution splat costs approximately 270\,MB per
scene---well within the storage economics of any institutional digital repository.

\begin{recommendation}
Adopt browser-based \acrshort{ksplat} delivery as the default public-access channel
for digitised cultural heritage reconstructions. Retain the full-resolution
\acrshort{usd} scene graph---including the uncompressed splat sublayer---as the
archival master. Derive compressed and mesh variants from this master on demand;
never treat a derivative as the sole preserved copy.
\end{recommendation}

% ─────────────────────────────────────────────────────────────────────────────
\section{Performance and Footprint}
\label{sec:delivery-performance}
% ─────────────────────────────────────────────────────────────────────────────

Table~\ref{tab:delivery-summary} summarises the key quantitative characteristics
of the delivery outputs for the gallery test scene.

\begin{table}[htbp]
  \centering
  \caption{Delivery output summary for the gallery test scene.}
  \label{tab:delivery-summary}
  \begin{tabular}{@{}L{5.5cm}R{4.5cm}@{}}
    \toprule
    \textbf{Metric}                           & \textbf{Value} \\
    \midrule
    Scene Gaussians (post-training)           & 1,180,000 \\
    Compressed \texttt{.ksplat} file size     & 18\,MB \\
    Browser render rate (1080p, mid-range \acrshort{gpu}) & 160\,FPS \\
    Browser render rate (1080p, integrated graphics) & 45--55\,FPS \\
    Environment mesh faces (\acrshort{tsdf} backend)  & $\approx$4.9\,M \\
    Environment mesh faces (\acrshort{milo} backend)  & $\approx$7.8\,M \\
    Per-object generative meshes              & 31 \\
    Archival \acrshort{usd} scene graph (with splat sublayer) & $\approx$270\,MB \\
    \midrule
    \multicolumn{2}{@{}l@{}}{\textit{Total pipeline wall-clock time (video to delivered scene)}} \\
    \quad\acrshort{tsdf} fast path            & $\approx$35\,min \\
    \quad \acrshort{milo} backend             & 1.5--2.5\,h (backend-dependent) \\
    \bottomrule
  \end{tabular}
\end{table}

The total wall-clock time from video ingest to a fully delivered scene---including
\acrshort{sfm}, Gaussian training, mesh extraction, per-object reconstruction, scene
graph assembly, and \acrshort{ksplat} compression---ranges from approximately
35~minutes on the \acrshort{tsdf} fast path to 1.5--2.5~hours when \acrshort{milo}
is selected as the extraction backend. The \acrshort{milo} range is explained by
the mesh-in-the-loop training overhead, which is sensitive to scene complexity and
the number of segmented objects.

For most cultural heritage digitisation workflows, where capture happens on a single
afternoon and processing overnight, even the upper bound of 2.5~hours represents a
negligible turnaround. The \acrshort{tsdf} fast path is appropriate when a same-day
preview delivery is required; the \acrshort{milo} path is preferred when the highest
mesh quality is needed for the engine delivery variant.

The \acrshort{gpu} processing throughout is dominated by \acrshort{3dgs} training
(approximately 23--28~minutes depending on variant, Section~\ref{sec:res-training}),
with mesh extraction and per-object reconstruction contributing the remaining time.
\acrshort{colmap} is the principal \acrshort{cpu}-bound stage at approximately
43~minutes, though it runs concurrently with early training initialisation in the
pipeline's asynchronous stage graph, reducing its contribution to the critical path.

\begin{keyfinding}
The three delivery paths---\acrshort{ksplat} browser streaming, \acrshort{usd}
engine import, and \acrshort{usd} archival master---are all produced from a single
processing run with no additional capture, no reprocessing, and no bespoke
per-channel optimisation. The incremental cost of supporting all three delivery
channels simultaneously is negligible relative to the cost of the reconstruction
itself.
\end{keyfinding}
